\lecture{7}{9. Marts 2026}{Volume and shape change}

\subsection{Compatibility}
Using $\epsilon_{ij} = 1 / 2 \left( u_{i,j} + u_{j,i} \right)$, the components of strain may be easily computed once displacements are known provided they are once differentiable.
However, the inverse problem of calculating displacements from strains is not as straight forward as there are six independent equations with only three unknowns.
Conditions of compatibility imposed on the components of strain $\epsilon_{ij}$, are both necessary and sufficient to ensure a continuous single-valued displacement field $\Vec{u}$.

The standard procedure is to eliminate the displacements between the equations repeatedly to produce equations with only strains as unknowns.
First we take two normal components $\epsilon_{( l l)}$ and $\epsilon_{(mm)}$ and the shearing component $\epsilon_{(lm)}$ with the summation convention suspended.

Taking second partial derivatives, we have
\begin{align*}
  \epsilon_{( l l,mm)} &= u_{(l,lmm)} \\
  \epsilon_{(mm, l l)} &= u_{(m,m l l)} \\
.\end{align*}
\begin{align*}
  \epsilon_{(lm,lm)} &= \frac{1}{2} \left( u_{(l,mlm)} + u_{(m,l lm)} \right) \\
  &= \frac{1}{2} \left( u_{(l,lmm)} + u_{(m,ml l)} \right)
.\end{align*}
Next we equate these to get
\begin{equation*}
  \epsilon_{(l l,mm)} + \epsilon_{(mm,l l)} = 2\epsilon_{(lm,lm)}
.\end{equation*}

Now, taking one normal component $\epsilon_{(l l)}$ and three shearing components $\epsilon_{lm}, \epsilon_{ln}$, and $\epsilon_{n n}$ and write the second partials
\begin{gather*}
  \epsilon_{(l l, mm)} = u_{(l,lmn} \\
  \epsilon_{(lm,l n)} = \frac{1}{2} \left( u_{(l,ml n)} + u_{(m,l l n )} \right) \\
  \epsilon_{(l n, lm)} = \frac{1}{2} \left( u_{(l, n lm)} + u_{(n,l l m)} \right)
  \epsilon_{(mn, l l)} = \frac{1}{2} \left( u_{(m, n l l)} + u_{(n, m l l)} \right)
.\end{gather*}
These can be equated to give
\begin{equation*}
  \epsilon_{(lm, l n)} + \epsilon_{(l n, lm)} - \epsilon_{(m n, l l)} = \epsilon_{(l l, mn)}
.\end{equation*}
These are six equations relating the components of strain that ensure that they will integrate into a unique displacement field.
These can be written explicitly as:
\begin{gather*}
  \epsilon_{11,22} + \epsilon_{22,11} = 2 \epsilon_{12,12} \\
  \epsilon_{22,33} + \epsilon_{33,22} = 2 \epsilon_{23,23} \\
  \epsilon_{33,11} + \epsilon_{11,33} = 2\epsilon_{31,31} \\
  \epsilon_{12,13} + \epsilon_{13,12} - \epsilon_{23,11} = \epsilon_{11,23} \\
  \epsilon_{23,21} + \epsilon_{21,23} - \epsilon_{31,22} = \epsilon_{22,31} \\
  \epsilon_{31,32} + \epsilon_{32,31} - \epsilon_{12,33} = \epsilon_{33,12}
.\end{gather*}
These are known as the \textit{St. Venant} compatibility equations and may also be extracted from the tensor equation
\begin{equation*}
  \epsilon_{ij,kl} + \epsilon_{kl,ij} = \epsilon_{ik,jk} + \epsilon_{jl,ik}
\end{equation*}
which represents $3^{4} = 81$ equations, where only six are distinct.
